\section{Introduction}\label{sec:intro}

\subsection{Two homologies for one knot}
Let $K\subset S^3$ be a knot. On the gauge-theoretic side sits the reduced singular instanton knot
homology $\Inat(K)$ of Kronheimer and Mrowka~\cite{KM1,KM2}. Its rank is bounded below by the
coefficient norm of the Alexander polynomial and above by the rank of reduced Khovanov homology
\eqref{eq:ell}. On the symplectic side sits the pillowcase program of Hedden, Herald and
Kirk~\cite{HHK1,HHK2}. Write $K=(Y,T)\cup(D,U)$ as the union of a tangle and the trivial tangle along a
Conway sphere. The traceless $\SU$ character variety of the four-punctured sphere is the
\emph{pillowcase}, a two-sphere with four orbifold points, and after a holonomy perturbation each side
determines an immersed curve in its smooth part: the tangle side a curve $L_1$, and the trivial side,
with an earring, the curve $L_0$, which depends on a small earring parameter $\epsilon>0$ (\S\ref{ss:earring}). Their Lagrangian Floer homology $\Hnat(Y,T,\pi)$ is defined when $L_1$ is restricted and the pair is admissible (Definitions~\ref{def:restricted} and~\ref{def:admissible}), and Hedden, Herald and Kirk conjecture that for suitable perturbations it is
isomorphic to $\Inat(K)$~\cite[Conj.~6.5]{HHK2} (Conjecture~\ref{conj:hhk65}), and they compare the $\Z/4$
gradings of the two sides~\cite[\S\S6.1, 11]{HHK2}.
For general decompositions figure-eight bubbles at the self-intersections of the curves deform the
differential. Cazassus, Herald, Kirk and Kotelskiy conjecture that each tangle is assigned an immersed
curve together with a bounding cochain, and that the deformed theory recovers
$\Inat(K)$~\cite[Conj.~D]{CHKK}; conditional on this conjecture, Smith proves that the cochain on one
Conway-sum tangle in $P(-2,3,5)$ cannot vanish~\cite[proof of Thm.~5.9]{Smith}.

This paper studies both sides for the pretzel knots $K_q=P(-2,3,q)$, $q\ge3$ odd (\S\ref{ss:knots}). The
smallest members are the torus knots $K_3=T(3,4)$ and $K_5=T(3,5)$, whose instanton and pillowcase
homology are treated in~\cite{WuebbenI,WuebbenIV,WuebbenV}, and for $q\ge7$ the knot $K_q$ is
hyperbolic. We decompose $K_q$ by regluing the tangle of Hedden, Herald and Kirk for $T(3,5)$ to the
trivial tangle by Dehn twists along the Conway sphere. For this decomposition we compute the pillowcase
homology with zero bounding cochain and find that it agrees with $\Inat(K_q)$; this proves the conjecture of
Hedden, Herald and Kirk for every $K_q$, and the same method applies to a family of twisted torus knots
(Theorems~\ref{thm:hhk} and~\ref{thm:twisted}). Smith decomposes the same knots as a sum of rational tangles,
with the earring on the $-2$ band. For that decomposition the pairing with zero bounding cochain has the wrong
rank at $q=5$~\cite[proof of Thm.~5.9]{Smith}, and also at $q=7$ for a piecewise-linear model of Smith's
curve~\cite{WuebbenIIb}. The companion paper~\cite{WuebbenIIb} constructs, for $q=5,7$, Maurer--Cartan elements
that correct the rank, in the two-arc algebra of Kotelskiy, Watson and Zibrowius~\cite{KWZ}; the instanton
dimensions of Theorem~\ref{thm:inat} are its targets.

\subsection{The instanton side}
\begin{theorem}\label{thm:inat}
For every odd $q\ge3$, $\Inat(K_q;\Z)\cong\Z^{q+2}$. For $q\ge7$ the signature is $\sigma(K_q)=-(q+1)$,
and in Kronheimer and Mrowka's absolute grading the graded ranks of $\Inat(K_q)$ are
\[
\bigl(1+N_3,\ N_1,\ N_1,\ N_3\bigr)\quad\text{in degrees }\gr-\sigma(K_q)=0,1,2,3,\qquad
N_1=\bigl\lceil\tfrac{q+1}4\bigr\rceil,\ N_3=\bigl\lfloor\tfrac{q+1}4\bigr\rfloor .
\]
\end{theorem}

The rank over $\Q$ is due to Lobb and Zentner~\cite[Prop.~4.7]{LZ}. For $q\ge7$ the freeness and the
graded ranks follow from the computation of the irreducible instanton homology of $K_q$ by Daemi and
Scaduto~\cite[Prop.~8.30]{DS24} together with their integral spectral sequence
(Proposition~\ref{prop:inat-graded}). Our proof of freeness is uniform in $q$: it runs Lobb and
Zentner's two-sided rank bound with integer coefficients, using Manion's integral computation of Khovanov
homology~\cite{Manion} and the integral spectral sequence of Kronheimer and Mrowka. For $q=3,5$ the
graded ranks are those of the torus knots in~\cite{WuebbenV}. Section~\ref{sec:inat} gives the proof and
Remark~\ref{rem:inat-attribution} the remaining attributions.

\subsection{The conjecture of Hedden, Herald and Kirk for the pretzel knots}
Hedden, Herald and Kirk compute the pillowcase homology of $T(3,5)$ from a tangle $(Y,T)$ whose
complement has free fundamental group of rank two, glued to the trivial tangle~\cite[\S\S10--11]{HHK2}.
Let $t_e$ be the Dehn twist along the curve of the Conway sphere that separates two of the punctures from
the other two. Regluing by $t_e^{-m}$ adds $m$ full twists on two strands, and the resulting knot is
$K_{5+2m}$ (Theorem~\ref{thm:topology}). That the original gluing gives the positive torus knot $T(3,5)$ is verified by a
computer-assisted argument (Remark~\ref{rem:chirality}). The unperturbed character variety $R(Y,T)$ is the union
of an arc and a circle that meet in two points $c_\pm$ (\S\ref{ss:straight}). The holonomy perturbations on the
tangle side are those of~\cite[\S10]{HHK2} used in~\cite{WuebbenV}: they are supported near two curves in $Y$,
with perturbation functions $\varepsilon_A\sin$ and $\varepsilon_B\sin$, and items (P1)--(P9) in
\S\S\ref{ss:straight}--\ref{ss:perturbed-P} recall the facts about them that we use from~\cite{WuebbenV}. For $a>0$ the closed cone
$\mathcal C_a$ is the set of parameters $\varepsilon=(\varepsilon_A,\varepsilon_B)$ such that, at both $c=c_+$ and $c=c_-$, the
linear part of $\kappa_c(\varepsilon)$, the critical value near $c$ of the function to which the equations of the
perturbed variety reduce near $c$, has absolute value at least $a|\varepsilon|$~\cite[Def.~5.5]{WuebbenV}; it is the
complement of an open conical neighborhood of the two lines $\{\varepsilon_B=\mp\sqrt3\,\varepsilon_A\}$
(Proposition~\ref{prop:restype}). The generator $r_+$ below is the intersection point of $L_0$ and $L_1$ near the
corner $(0,0)$ that comes from the abelian representation of the knot group (\S\ref{ss:decomp}).

\begin{theorem}\label{thm:hhk}
Let $q\ge7$ be odd, and let $K_q=(Y,T)\cup_{t_e^{-(q-5)/2}}(D,U)$ be the decomposition above. Fix $a>0$.
For every sufficiently small $\varepsilon\in\mathcal C_a\setminus\{0\}$ and every sufficiently small earring
parameter:
\begin{enumerate}
\item $L_1$ is a restricted immersed $1$-manifold, and $(L_0,L_1)$ is admissible;
\item the Floer complex has $q+2+4N_q$ generators, where $N_q=\lfloor 5(q-6)/12\rfloor-\lfloor(q-6)/12\rfloor$;
\item the differential has rank $2N_q$, and it is nonzero for $q\ge9$;
\item with the absolute grading normalized by $\gr(r_+)=\sigma(K_q)$, as in~\cite[(19)]{HHK2}, and with any absolute grading on the circle component of $L_1$ when there is one, there is an isomorphism $\Hnat(Y,T,\pi)\cong\Inat(K_q;\F_2)$ of $\Z/4$-graded vector spaces.
\end{enumerate}
In particular, the conjecture of Hedden, Herald and Kirk~\cite[Conj.~6.5]{HHK2} holds for every $K_q$.
\end{theorem}

The perturbation resolves the two double points $c_\pm$ of $R(Y,T)$ in one of two ways, joined for
$|\varepsilon_B|<\sqrt3\,|\varepsilon_A|$ and split for $|\varepsilon_B|>\sqrt3\,|\varepsilon_A|$ (\S\ref{ss:perturbed-P} and
Proposition~\ref{prop:restype}); for small $a$ both occur in $\mathcal C_a$, and the theorem holds in both. For $q\ge9$ the complex is not minimal. Hedden, Herald and
Kirk verify the conclusion of their conjecture for two-bridge knots~\cite[\S7]{HHK2}, a family that
contains hyperbolic knots, and for several torus knots, among them $T(3,4)$, $T(3,5)$ and $T(3,7)$~\cite[\S\S10--11]{HHK2}; the torus knots
$T(3,n)$ are treated in~\cite{WuebbenV}. Fukumoto, Kirk and Pinz\'on-Caicedo describe the unperturbed
character variety of the pretzel knots for a decomposition of their own and count $9$ generators for
$P(-2,3,7)$~\cite[\S5]{FKP}
(Remark~\ref{rem:FKPcount}). For $q\ge7$ the knot $K_q$ is hyperbolic and not two-bridge
(Proposition~\ref{prop:hyperbolic}). To our knowledge, these are the first hyperbolic knots outside the
two-bridge family for which the conjecture is verified, with the $\Z/4$ gradings and nonzero Floer
differentials.

The same regluing applies to the tangles of~\cite[\S10]{HHK2} for $T(3,n)$. The \emph{twisted torus knot}
$T(3,n;2,m)$ is the closure of the braid $(\sigma_1\sigma_2)^n\sigma_1^{2m}$.

\begin{theorem}\label{thm:twisted}
Let $n\ge7$ with $n\equiv1$ or $2\pmod 6$, let $m\ge1$, and let $K=T(3,n;2,m)$. Then
$\sigma(K)=\sigma(T(3,n))-2m$, and $\Inat(K;\Z)$ is free of rank $|\sigma(K)|+1$. Assume, as in~\cite{HHK2}
and~\cite{WuebbenV}, that the tangle of~\cite[\S10]{HHK2} for $T(3,n)$ glued to the trivial tangle is the positive
torus knot $T(3,n)$. Then this tangle, reglued by $t_e^{-m}$, is a decomposition of $K$, and for small holonomy
perturbations and generic small earring parameters it satisfies the conclusion of~\cite[Conj.~6.5]{HHK2}:
$\Hnat\cong\Inat(K;\F_2)$ as $\Z/4$-graded vector spaces.
\end{theorem}

Theorem~\ref{thm:H} gives the precise form, with the graded ranks. The family of
Theorem~\ref{thm:twisted} contains torus knots, $T(3,6j+1;2,1)=T(3,6j+2)$, and by braid conjugacy
$T(3,6j+2;2,m)=T(3,6j+1;2,m+1)$ (Lemma~\ref{lem:conj}). The other members, $T(3,6j+1;2,M)$ with $j\ge1$
and $M\ge2$, are not torus knots, not alternating and not pretzel knots
$P(-2,3,q)$ (Lemma~\ref{lem:twisted-not-torus}). We make no claim about their hyperbolicity.

\subsection{Outline of the arguments}
\emph{Theorem~\ref{thm:hhk}.} The regluing replaces $L_1$ by its image under the linear shear
$(\gamma,\theta)\mapsto(\gamma,\theta+(q-5)\gamma)$ of the pillowcase (Lemma~\ref{lem:shear}). The unperturbed
character variety of the $T(3,5)$ tangle is the union of an arc of binary dihedral characters and a circle of characters, the \emph{central circle}, meeting it in two
points. A word $w$ in the tangle group with $\rho(w)=1$ for every representation $\rho$ on the central
circle shows that its image in the pillowcase is a union of straight segments (Lemma~\ref{lem:straight} and
Corollary~\ref{cor:polygon}). The perturbed curves are then those of~\cite{WuebbenV}, sheared, and on them
we compute the intersection points with $L_0$, their $\Z/4$ gradings, and the bigons by hand. The
bigons come in chains along the lifts of the earring curve $L_0$ to $\R^2$, which gives the rank $2N_q$ of the
differential (Lemma~\ref{lem:zigzag}). The identification of the reglued knot uses only its
fundamental group: the relator is that of Clay and Watson's presentation of the twisted torus
knots~\cite[Prop.~22]{ClayWatson} (Lemma~\ref{lem:relator}), two-generator knots are prime~\cite{Norwood},
knots are determined by their complements~\cite{GordonLuecke}, and an induction on the signature fixes the
chirality.

\emph{Theorem~\ref{thm:twisted}.} For $n\equiv1,2\pmod6$ the unperturbed variety is an arc with a straight
image and $2\lfloor n/6\rfloor$ circles that miss the edges of the pillowcase~\cite{WuebbenV}. The sheared
arc meets $L_0$ in $1+2M$ points and carries no bigon, and each circle contributes one dimension in each
degree to the homology, by a theorem of Hedden, Herald and Kirk~\cite[Thm.~5.4]{HHK2}. On the knot side, a closed formula for
the Burau matrix of the braid gives the signature and the coefficient norm of the Alexander polynomial,
and the instanton rank follows from the work of Daemi and Scaduto.

\subsection{Scope of claims}
Theorem~\ref{thm:inat} is unconditional. Theorems~\ref{thm:hhk} and~\ref{thm:twisted} rest on the analysis of
the perturbed curves of the torus-knot tangles in~\cite{WuebbenV}. For Theorem~\ref{thm:hhk} the base case is verified by
computer (Remark~\ref{rem:chirality}); for the decomposition in Theorem~\ref{thm:twisted} it is a hypothesis of the
statement, shared with~\cite{HHK2} and~\cite{WuebbenV}, that the decomposition of $T(3,n)$ gives the positive torus
knot, which the knot group cannot detect. Theorem~\ref{thm:twisted} also uses the
computation of Hedden, Herald and Kirk for circles that miss the edges~\cite[Thm.~5.4]{HHK2}, whose proof
in~\cite{HHK2} rests on their invariance theorem~\cite[Thm.~4.1]{HHK2} (Remark~\ref{rem:invariance}). Numerical traces of the perturbed
curves agree with them (Remark~\ref{rem:numericsP}) but are not used in the proofs.

\subsection{Organization}
Section~\ref{sec:background} recalls the pillowcase Floer theory of Hedden, Herald and Kirk.
Section~\ref{sec:inat} proves Theorem~\ref{thm:inat}. Section~\ref{sec:sheared} proves Theorem~\ref{thm:hhk},
and Section~\ref{sec:twisted} proves Theorem~\ref{thm:twisted}. Section~\ref{sec:disc} discusses related work
and open problems, and Appendix~\ref{app:data} describes the verified computation behind
Remark~\ref{rem:chirality} and the computations that corroborate the proofs.

\subsection*{Acknowledgements}
The pillowcase model and the geometry used here are due to Hedden, Herald and Kirk~\cite{HHK1,HHK2} and to
Fukumoto, Kirk and Pinz\'on-Caicedo~\cite{FKP}. The closed-form Khovanov input is Manion's~\cite{Manion}, the Alexander polynomials belong to Hironaka's family~\cite{Hironaka}, the
presentation of the twisted torus knot groups is Clay and Watson's~\cite{ClayWatson}, the rational rank of
the instanton homology is Lobb and Zentner's~\cite{LZ}, and the integral and graded instanton computations
for $q\ge7$ and for twisted torus knots also follow from the work of Daemi and Scaduto~\cite{DS24}.
